Feature-composed I-ProtoMF, fI-ProtoMF, changes the host model in exactly
one place. I-ProtoMF~\cite{melchiorre2022protomf} represents each item by
a free embedding, one learned row per item in a lookup table. fI-ProtoMF
replaces this table by a composition. The embedding of item $i$ is the
sum of the embeddings of its recorded feature values,
\begin{equation}
\mathbf{q}_i \,=\, \sum_{f \in \mathcal{F}_i} \mathbf{e}_f
\,+\, \mathbf{e}_{\mathrm{ID}(i)},
\label{eq:fi-composition}
\end{equation}
where $\mathcal{F}_i$ holds the feature values of item $i$, each
$\mathbf{e}_f \in \mathbb{R}^d$ is a learned vector shared by all items
that carry the value $f$, and $\mathbf{e}_{\mathrm{ID}(i)}$ is an
optional per-item row whose role is discussed below. The composition
follows the established template of feature-aware matrix
factorization~\cite{chen2012svdfeature,kula2015lightfm}, applied here to
the item branch of a prototype model. Everything else remains the host.
The item prototypes $\mathbf{p}_1, \dots, \mathbf{p}_K$ stay free
learned vectors, the activation of prototype $k$ is still the shifted
cosine similarity $t^*_{i,k} = 1 + \cos(\mathbf{q}_i, \mathbf{p}_k)$,
the user is still a free vector $\mathbf{u} \in \mathbb{R}^K$ over
prototype activations, and the score
$S(u,i) = \mathbf{u} \cdot \mathbf{t}^*_i$ is trained with the host's
loss and prototype regularizers (Section~\ref{sec:bg-protomf}).

Two terms carry this chapter and are fixed here. An attribute is a
human-readable catalogue property, the department of a garment or its
colour group. A feature is the encoded, model-facing form of such a
property, in this model a learned row $\mathbf{e}_f$ entering the sum in
Equation~\ref{eq:fi-composition}. Statements about meaning speak of
attributes, statements about mechanism speak of features. The model
names keep the f of feature rather than attribute because the
composition sums whatever encoded rows it is given, an ID handle or an
engineered row as readily as an attribute value, and choosing
attribute-derived rows is what makes the resulting model interpretable.

The ID row makes the host a special case of the composition. When the
feature set $\mathcal{F}_i$ is empty and the ID row is kept,
$\mathbf{q}_i = \mathbf{e}_{\mathrm{ID}(i)}$ is again a free per-item
vector, and the model reduces to I-ProtoMF exactly. Our implementation
reproduces the host bit-identically in this configuration, so every
difference observed between fI-ProtoMF and its host is caused by the
composed embedding and by nothing else. Switching the two ingredients of
the sum yields the three arms evaluated in this chapter. The headline
arm, fI-ProtoMF, uses features and the ID row. The noid arm drops the ID
row and represents every item from features alone, which isolates what
the recorded vocabulary supports by itself. The f0 arm drops the
features and keeps the ID row. It is the host reached through the
extended code path and serves as the reference point of all comparisons.

The change redraws which components of the model carry meaning. The item
embeddings are grounded by construction. Every $\mathbf{q}_i$ is a sum
of rows, and each row stands for a catalogue property or for the item's
identity, so the make-up of an item representation can be stated in
words. The prototypes remain free learned vectors. They are not grounded
themselves, but they now anchor a space populated by attribute rows, and
this shared residence is what later allows their meaning to be read from
parameters alone (Section~\ref{sec:fi-readouts}). The user vector
remains a purely collaborative object with no vocabulary attached. This
division draws a claim boundary the thesis observes throughout.
Explanations of this model are stated in attribute vocabulary, because
the item side of the score is built from attribute rows. The weights
that combine the activations into a score stay learned collaborative
quantities, so the model's reasoning is never claimed to be
feature-transparent. On this host the boundary is clean. The activation
vector $\mathbf{t}^*_i$ carries the entire dependence of the score on
the item, so the grounded vocabulary covers everything the model states
about items, while everything it states about the user stays latent by
design.

The composition also determines what the model can do with an item it
has never seen. Because the score is bilinear, one informative side
suffices for a meaningful prediction. A cold item has no trained ID row,
but its feature values are known the moment it enters the catalogue, so
Equation~\ref{eq:fi-composition} yields a usable $\mathbf{q}_i$ from the
feature rows alone. At cold scoring the untrained ID row is dropped from
the sum. Cold users are the mirrored case and stay structurally out of
reach for an item-side design, since helping them requires a vocabulary
on the user side. Cold-start claims for this model are therefore scoped
to cold items met by warm users, a scope that matches the catalogue
churn which makes the fashion domain item-cold in the first place. The
user side receives its own treatment in Chapter~\ref{ch:fu}.
